\subsection{多固有値領域の三次共鳴による必要被覆}
\label{subsec:cubic-spectral-resonance}
\begin{proposition}[単純固有直線の孤立純三乗共鳴の排除]
\label{prop:cubic-spectral-resonance}
H3包絡の中山自己同型の半単純部分について
$B_1=\bigoplus_{i=1}^r V_i$、固有値を相異なる $\lambda_i$、
$m_i=\dim V_i$ とする。ある非負整数組 $a=(a_1,\ldots,a_r)$ が
\[
 |a|=3,\qquad \lambda^a=1
\]
を満たすことが必要である。
さらに、そのような組がただ一つで $a=3e_i,m_i=1$ ならH3は存在しない。
これは中山の単冪部分が非自明でも成立する。
\end{proposition}
\begin{proof}
定理\ref{thm:nakayama-cubic-linearization}より $\mu_B(G)=G\ne0$。
有限次元空間 $B_3$ 上のJordan分解から $\mu_{B,s}(G)=G$。
積写像 $B_1^{\otimes3}\twoheadrightarrow B_3$ は半単純作用に関して同変であり、
固有値1の部分上も全射である。tensor側の固有値は、各語の文字数 $a$ に対する
$\lambda^a$ なので最初の主張を得る。
唯一の共鳴が $3e_i$ で $V_i=kx$ なら、固有値1部分の像は $kx^3$ に含まれる。
従って $G=cx^3$、$c\ne0$ である。
$R=B/(G)$ では $x$ と $x^2$ はそれぞれ次数1、2の非零元のまま残る一方、
$x\,x^2=0$。これはH3の整域性に反する。
\end{proof}

\paragraph{固有値重複型 $(2,1,1)$ の具体的な被覆．}
重複2の固有値を $a$、単純な二つを $b,c$ とする。
必要共鳴は
\[
 a^3,a^2b,a^2c,ab^2,abc,ac^2,b^3,b^2c,bc^2,c^3
\]
のいずれかが1となることである。
ただ一つの共鳴が $b^3=1$ または $c^3=1$ の領域は空である。
正確な被覆は、八共鳴の和集合と $\{b^3=c^3=1\}$ の和集合である。
この最後の有限固有値条件を省略しない。

\paragraph{四単純固有値．}
$\lambda_1,\ldots,\lambda_4$ が相異なる場合、20個の三次指数組がある。
四つの純三乗共鳴それぞれについて、それ以外の19共鳴が成立しない開部分は空である。
従って候補は16個の非純共鳴の和集合と、少なくとも二つの純共鳴が成立する領域で覆われる。
これらは必要条件のみであり、AS4存在、正常核の非消滅、H3十分性を含まない。
特に固有値の重複や複数共鳴の交叉を一般点の排除から消去していない。

\begin{theorem}[四単純固有値の二次非共鳴領域はskew]
\label{thm:cubic-simple-spectrum-skew}
$B$ を上記のAS4包絡とし、中山の半単純部分が $B_1$ 上に相異なる四固有値
$\lambda_1,\ldots,\lambda_4$ を持つとする。
十個の積 $\lambda_i\lambda_j$ ($i\le j$) が全て相異なるなら、
$B$ はその固有基底に関してskew polynomial algebraである。
従って未分類の非skew四単純固有値候補は次の21条件のいずれかを満たす：
\[
 \lambda_i=-\lambda_j\ (i<j),\qquad
 \lambda_i^2=\lambda_j\lambda_k\ (i,j,k\text{相異なる}),\qquad
 \lambda_i\lambda_j=\lambda_k\lambda_l\ (\{i,j,k,l\}=\{1,2,3,4\}).
\]
中央の条件は $i$ と順序なしの $\{j,k\}$ で12個、最後は分割として3個と数える。
\end{theorem}
\begin{proof}
二次関係空間を $I_2\subset B_1^{\otimes2}$ とする。
その次元は6であり、中山の半単純作用に不変である。
十個の固有値が異なるので、tensorの固有空間は
$kx_i^2$ と $\langle x_ix_j,x_jx_i\rangle$ ($i<j$) である。
整域性から $I_2\cap kx_i^2=0$。
二次元の後者との交叉は高々一次元で、非零なら両係数とも非零である。
一方の係数が零なら非零一次元の積が零になるからである。
六つのpair空間しかなく $\dim I_2=6$ なので、全てのpairに関係が一本ずつ存在する。
従って $I_2$ は $x_ix_j-q_{ij}x_jx_i$、$q_{ij}\ne0$ の六本で張られる。
$B$ はquadraticなので、これが全提示である。
最後に二つの二次指数組が添字を共有すれば、その積の等式は異なる固有値の等式に帰着して不可能。
共有しない場合は、平方同士、平方と二つの異なる添字、二つの互いに素なpairという
三種類だけであり、列挙した6+12+3条件を得る。
\end{proof}

\begin{theorem}[重複型 $(2,1,1)$ の二次非共鳴領域]
\label{thm:cubic-211-nonresonant-Ore}
中山の半単純部分の固有空間を $U\oplus kb\oplus kc$、
固有値を $a,b_0,c_0$ とし、$\dim U=2$ とする。
六個の積 $a^2,ab_0,ac_0,b_0^2,b_0c_0,c_0^2$ が全て相異なるなら、
包絡は二元AS2平面 $T=T(U)/(r)$ 上の逐次skew polynomial拡大
\[
 B=T[b;f][c;\widetilde g],\qquad
 f,g\in\operatorname{Aut}_{\rm gr}T,\quad fg=gf,\quad
 \widetilde g|_T=g,\quad\widetilde g(b)=q^{-1}b
\]
となる。ここで $q\ne0$。中山の単冪部分の消滅は仮定していない。
\end{theorem}
\begin{proof}
二次tensorの六固有空間は $U^{\otimes2}$、$Ub+bU$、$Uc+cU$、
$kb^2$、$kbc+kcb$、$kc^2$ である。
整域性から関係空間はdecomposableな非零tensorを含まない。
$U^{\otimes2}$ と関係空間との交叉は高々一次元である。
実際、二次元の $2\times2$ 行列pencilは代数閉体上で非零特異行列を含む。
$Ub+bU$ との交叉は両射影が単射なので高々二次元、$Uc+cU$ も同様。
$b^2,c^2$ は関係に入らず、$bc,cb$ の空間には高々一本である。
上界の和は $1+2+2+1=6$ であり、関係数が6なので全て等号となる。
従って関係は、非退化 $r\in U^{\otimes2}$、
$b u=f(u)b$、$c u=g(u)c$ ($u\in U$)、$bc=qcb$ と書ける。
$f,g$ は可逆で $q\ne0$。

二元非退化二次関係は二元AS2平面を定める。
$br=f(r)b=0$ と整域性から $f(r)=0$ in $B_2$ であり、
$B_2$ の純 $U$ 関係が $kr$ なので $f(r)\in kr$。
従って $f$ は $T$ の自己同型であり、$g$ も同様。
$(bc)u$ を二通りに書き換えて差を取ると
$(fg-gf)(u)bc=0$。整域性から $fg=gf$。
従って表示した $\widetilde g$ は $T[b;f]$ の自己同型を定める。
得た逐次Ore拡大と $B$ は同じ二次生成関係を持ち、$B$ はquadraticなので一致する。
\end{proof}

\subsubsection{三固有値共鳴の四つの具体的chart}
\begin{proposition}[二次・三次共鳴の有限スペクトルatlas]
\label{prop:cubic-three-spectrum-atlas}
重複2の固有値を先頭に置き、二つの単純固有値の交換を許す。
系\ref{cor:cubic-multispectrum-quadratic-resonance}の三固有値必要被覆は、次の四chartで覆われる。
全パラメータは非零とし、表示した固有値は相異なるものに制限する。
\[
\begin{array}{c|c|l}
 & (a,b,c)&\text{必要な三次共鳴（少なくとも一つ）}\\\hline
 A&(a,-a,c)&a^3=\pm1,\ a^2c=\pm1,\ ac^2=\pm1,\ c^3=1\\
 B&(a,b,-b)&a^3=1,\ a^2b=\pm1,\ ab^2=\pm1,\ b^3=\pm1\\
 C&(a,at,a/t)&a^3t^k=1\quad(-3\le k\le3)\\
 D&(a,at,at^2)&a^3t^k=1\quad(0\le k\le6)
\end{array}
\]
$C,D$ では $t\ne\pm1$。これは必要スペクトルのatlasであり、包絡やH3対象の存在を保証しない。
さらに $C$ の $k=\pm3$ 又は $D$ の $k=6$ の枝がH3に現れるなら、
$t$ の位数は $3,4,5,6$ のいずれかである。
$A$ の $c^3=1$ 枝がH3に現れるなら、$a^{12}=1$ が必要である。
従ってこれらの純共鳴枝の非捩れ一般点を候補として保持する必要はない。
\end{proposition}
\begin{proof}
六つの二次共鳴は、単純固有値の交換を除くと
$a=-b$、$b=-c$、$a^2=bc$、$b^2=ac$ の四場合である。
前二つは $A,B$、後二つは $t=b/a$ と置いて $C,D$ の表示を得る。
十個の三次指数組に代入して同じ単項式をまとめると、各行の七条件を得る。

$C$ の $k=3,-3$ はそれぞれ単純固有直線 $b,c$ の純三乗だけに対応する。
$D$ の $k=6$ も $c$ の純三乗だけである。
孤立純共鳴の排除により、これらの枝では異なる $k'$ の共鳴がもう一つ必要。
従って $t^{k-k'}=1$、$1\le|k-k'|\le6$。
$t\ne\pm1$ を合わせると位数は3,4,5,6に限られる。
$A$ の $c^3=1$ でも、別の共鳴が必要である。
$a^3=\pm1$、$a^2c=\pm1$、$ac^2=\pm1$ のいずれも、$c^3=1$ と合わせれば $a^{12}=1$ を含意する。
有限捩れ点の全てがH3であるという逆向きは主張しない。
\end{proof}

\begin{corollary}[半単純固有値集合の次元上界]
\label{cor:cubic-spectral-dimension}
三固有値以上のH3包絡が取り得る半単純固有値の集合のZariski閉包は、
三固有値の場合は次元高々1、四固有値の場合は次元高々2である。
包絡族や内在的moduli自体の次元の上界ではない。
\end{corollary}
\begin{proof}
$r$ 個の相異なる固有値のtorus上で、必要な二次共鳴は
$\lambda^u=\epsilon$、$\sum u_i=0$、$u\ne0$、$\epsilon\in\{1,-1\}$、
三次共鳴は $\lambda^v=1$、$\sum v_i=3$ の形となる。
$u,v$ は有理数上独立なので、両条件の共通部分は空か、
次元 $r-2$ のtorusの有限個のtranslateである。
有限個の共鳴の和集合と固有値の有限置換で次元は増えない。
\end{proof}

\begin{lemma}[四次topによる追加スペクトル条件]
\label{lem:cubic-quartic-spectral-support}
H3包絡の相異なる中山半単純固有値を $\lambda_i$ とする。
四次指数組 $v$ で $|v|=4,\lambda^v=1$ となるものが存在する。
さらに各固有空間 $V_i$ に対し、$v_i>0$ のそのような組が存在する。
単冪部分の消滅は不要である。
\end{lemma}
\begin{proof}
AS4型14641のKoszul双対Frobenius代数の最高tensorを $w\in B_1^{\otimes4}$ とする。
本稿の規約では $\nu=\mu_B^{-1}$、右巡回を $\tau$ として
$(\nu\otimes1^{\otimes3})\tau(w)=-w$ である。
四回反復して $\nu^{\otimes4}w=w$、従って $\mu_B^{\otimes4}w=w$ を得る。
半単純部分についても同じ不変性が成立する。
$w$ の各非零語の文字数が必要な $v$ を与える。
Frobeniusの次数1と3のperfect pairingにより、
$B_1^*\to B_1^{\otimes3}$ という $w$ の縮約は単射である。
従って $w$ はどの非零固有空間の文字も全く含まないということはない。
\end{proof}

\begin{theorem}[三固有値スペクトルの有限捩れ部分と二つの残余pencil]
\label{thm:cubic-three-spectrum-torsion-pencil}
三つの相異なる中山半単純固有値を持つH3包絡では、重複2の固有値を先頭に書くと、
次の少なくとも一つが成立する。
\begin{enumerate}
\item ある $1\le N\le24$ に対して全固有値が $N$ 乗根である。
\item $(a,b,c)=(1,z,z^{-1})$、$z\ne\pm1$。
\item $(a,b,c)=(z,1,z^{-1})$、$z\ne\pm1$。
\end{enumerate}
二つの単純固有値の交換は許す。
後二つのpencilや有限捩れ点のH3十分性は主張しない。
\end{theorem}
\begin{proof}
命題\ref{prop:cubic-three-spectrum-atlas}の四chartを使う。
$A$ では三次共鳴は $a^ic^{3-i}=\epsilon$、四次共鳴は
$a^rc^{4-r}=\epsilon'$、$\epsilon,\epsilon'\in\{1,-1\}$ の形となる。
指数行列の行列式は $d=4i-3r$。
$i=1,2$ なら零にならない。
$i=0$ なら前補題で $r>0$ の四次語を選べ、$i=3$ なら $r<4$ の四次語を選べるため、
いずれも $d\ne0$ とできる。
従って $a^{2|d|}=c^{2|d|}=1$ であり、$2|d|\le24$。
$-a$ も同じ偶数乗で1となる。
$B$ も $a,b$ の指数と符号を用いて全く同様である。

$C$ では三次 $a^3t^k=1$、四次 $a^4t^l=1$、
$-3\le k\le3,-4\le l\le4$ であり、行列式は $d=3l-4k$。
どれかの非零四次語で $d\ne0$ なら $a^{|d|}=t^{|d|}=1$、$|d|\le24$ となる。
そうでなければ $k=0,3,-3$ であり、それぞれ $l=0,4,-4$ しか非零四次語に現れない。
三次・四次式の最大公約数 $\gcd(3,4)=1$ から
それぞれ $a=1,at=1,a/t=1$ を得る。
$k=0$ は第二pencilである。
$k=\pm3$ では全非零四次語が単純固有直線の純四乗となり、前補題の非退化性に反する。

$D$ では $0\le k\le6,0\le l\le8$ で同じ行列式を使う。
$d\ne0$ なら同様に $|d|\le24$ の第一項となる。
残りは $(k,l)=(0,0),(3,4),(6,8)$。
$(0,0)$ では $a=1$ で、全非零四次語が重複固有空間の文字だけになる。
これは前補題の非退化性に反する。
$(3,4)$ では $at=1$ となり、$z=a$ と置けば第三pencilである。
$(6,8)$ では全非零四次語が最後の固有直線の純四乗となり、同じく非退化性に反する。
以上は全て必要条件の議論であり、有限捩れ点を無条件に排除していない。
\end{proof}
